\documentclass[12pt,a4paper]{article}

% ------------------------------------------------------------
% Basic packages
% ------------------------------------------------------------
\usepackage[utf8]{inputenc}
\usepackage[T1]{fontenc}
\usepackage[english]{babel}

% Page layout
\usepackage{geometry}
\geometry{margin=1in}

% Math
\usepackage{amsmath}
\usepackage{amssymb}
\usepackage{mathtools}

% Tables and figures
\usepackage{graphicx}
\usepackage{booktabs}
\usepackage{adjustbox}
\usepackage{float}

% References and links
\usepackage[hidelinks]{hyperref}

% Optional: better spacing
\usepackage{setspace}
\onehalfspacing

% ------------------------------------------------------------
% Title information
% ------------------------------------------------------------
\title{Home Bias and International Risk Sharing}
\author{Your Name}
\date{\today}

% ------------------------------------------------------------
% Document starts here
% ------------------------------------------------------------
\begin{document}

\maketitle

\begin{abstract}
This paper investigates the relationship between equity home bias and international risk sharing.
The analysis is based on panel-data regressions using OECD national accounts data and foreign asset holdings.
\end{abstract}

\newpage

\tableofcontents

\newpage

% ============================================================
\section{Introduction}
% ============================================================
In their 2007 Paper Sørensen et al. argue that risk sharing and home bias, are two economic patterns 
that were analysed seperately before but "may be manifestations of the same underlying
behavior".

After replicating the main results of the paper, and extending it by time and sample size\footnote{As much as possible,
constrained by data availability.} we shift our focus the eurozone as an interesting case study, to further investigate how 
home bias and international risk sharing develops. European Integration during the original papers sample period (1993--2003) 
and afterwards presents a fascinating opportunity to investigate, whether a decline in the exchange rate risk and 
transaction costs, two main reasons for Home Bias, is actually associated with a declining Home Bias and increasing 
risk sharing. Naive data analysis finds that the Euroarea countries became more diversified during 
their integration, but not within themselves, but rather with countries outside the Eurozone. This highly puzzling result
motivated us to analyse the subject further, finding that Financial Centers inside Europe are the main drivers of this trend 
and present a major unsolved challange to the original Sorensen et al. paper. Hence, rather than doing a single case study, we 
shifted our papers focus to analyse this shortcoming more throughoutly.

% ============================================================
\section{Panel regressions}
% ============================================================
\subsection{Theoretical outline}

The outhors use a panel regression setting to analyse
their central question - whether a lower home bias is associated
with a higher level of income risk sharing and consumption risk sharing. 
To do so, they regress the log growth rates of country specific Gross National 
Income (GNI) or Consumption (C) relative to the full sample's log growth rate of GNI 
or Consumption on the country specific log Gross National Product (GDP) growth rate, 
relative to the full samples GDP growth rate.  
 The baseline regressions are

\begin{equation}
\Delta \log GNI_{it} - \Delta \log GNI_t
=
\alpha_i
+
\kappa
\left(
\Delta \log GDP_{it} - \Delta \log GDP_t
\right)
+
\epsilon_{it}
\end{equation}

and 
\begin{equation}
\Delta \log C_{it} - \Delta \log C_t
=
\alpha_i
+
\eta
\left(
\Delta \log GDP_{it} - \Delta \log GDP_t
\right)
+
\epsilon_{it}.
\end{equation}

As in the cross sectional regressions,
one minus the respective risk sharing coefficents \(\kappa\) and \(\eta\) measures the 
average amount of risk sharing of a country.\\
To include the equity home bias into this regression, the authors apply
an interaction term for the equity home bias of a country, relative to 
the average equity home bias in the sample and an interaction term for a 
potential time trend. 

\begin{equation}
\kappa
=
\kappa_0
+
\kappa_1(t - \overline{t})
+
\kappa_2(EHB_{it} - \overline{EHB}_t)
\end{equation}

and 

\begin{equation}
\eta
=
\eta_0
+
\eta_1(t - \overline{t})
+
\eta_2(EHB_{it} - \overline{EHB}_t).
\end{equation} 

Once again it is most intuitive to interpret the average risk sharing 
coefficient of a country: 
\begin{equation}
AverageIncomeRiskSharing
=
1 -\kappa_0
-
\kappa_1(t - \overline{t})
-
\kappa_2(EHB_{it} - \overline{EHB}_t)
\end{equation}

and 

\begin{equation}
AverageConsumptionRiskSharing
=
1 - \eta_0
-
\eta_1(t - \overline{t})
-
\eta_2(EHB_{it} - \overline{EHB}_t).
\end{equation} 
For the hypothesis of a potential co-movement between
the equity home bias and risk sharing we must focus on the third coefficients \(-\kappa_2\) and \(- \eta_2\), 
implying that an increase in the equity home bias by 0.1 increases income risk sharing by \(-\kappa_2\)
and consumption risk sharing by \(- \eta_2\).


All panel regressions are run as two-stage feasable Ordinary Least Squares Estimations.
In the second stage regression each country is attributed
a weight relative to 1 over its error terms standard deviation in the
first stage regression.\footnote{Since this is weakly documented in the original paper, 
we compared the first and second stage regression results and conclude that the differences have no impact on the 
direction and significance of coefficients.}

\subsection{Empirical results}
When reporting the results of the panel regressions, we follow
the authors standard and report the easier to interpret 
coefficients \(1-\kappa_0/1-\eta_0\), \(-\kappa_1/-\eta_1\) and \(-\kappa_2/-\eta_2\) multiplied
by 100 to interpret them as percentage effects. Reassuringly, our panel regression results 
mainly follow the key results of the Sorensen et. al (HIER LINK EINFÜGEN) paper, with a statistically 
non-different to zero estimate for the time trend and a highly significant coefficient
for the equity home bias on both income and consumption risk sharing. An increase in the equity home bias by 0.1 
is now associated with an increase in income risk sharing by roughly 13 percent and consumption risk sharing by 
roughly 20 percent. Given that our point estimates for the coefficient are larger than 100 and 
notably higher than the point estimates found
in the Sorensen et al (HIER LINK EINFÜGEN) paper, we believe that data revisions and the choice of the 
base year (AS EXPLAINED ABOVE) play an important role. Further, we rely on an interpretation
offered by the authors themseves, believing that: "a large decline in home bias is not 
likely to lead to more than 100 percent risk sharing [...] and we interpret the coefficient to imply that declining home bias has been associated
with strongly increasing [...] risk sharing, even if the actual value is likely to not be valid for very large changes in home bias."
Further, our coefficient for the average income risk sharing is highly significant, while it is not 
in the original paper, something that we blame on data revisions and the choice of the base year. \\

In table 5, we follow the papers approach and run similar income risk sharing and consumption 
risk sharing regressions on terms of certain foreign asset holdings over GDP.\footnote{Note that a positive 
sign of \(-\kappa_2/ -\eta_2\) is now associated with a higher level of risk sharing.} 
We treat these regressions merely as an robustness check, and find postive \(-\kappa_2, -\eta_2\)coefficients 
for all forms of asset holdings, and statistically non signfiicant estimates for the time trend.\\

Even though we find partly differing point estimates, due to data revisions or the choice of the different base year, 
our findings is mainly in line with the authors finding for the same country and time sample. 
Therefore, we conclude that the main findings of the Sorensen et al (HIER PAPER EINFÜGEN) panel regression can be replicated. 

% ============================================================
\section{Extentions}
% ============================================================

% ============================================================
\subsection{First Extention: Similar Country Sample 1987--2017}
% ============================================================
We aim to analyse how the 22 countries of the original analysis have developed during a by 
20 years extended time period from 1987 to 2017. This extention is particularly interesting
as the original period of analysis of 11 years is rather short. 

We find that the results of the authors analysis and our replication largely hold when extending the 
time period by 20 years. A 0.1 decrease of the equity home bias is now associated with a roughtly 3 percent increase 
in income risk sharing (HIER TABELLE VERLINKEN) and a 8 percent increase in consumption risk sharing. 

The estimate for the time trend remains statistically non-significant and the average risk sharing coefficient 
for income and consumption risk sharing decreases notably with the coefficient for average income risk sharing becoming
statistically non significant whilst the coefficient for average consumption risk sharing remains highly significant. 
Given that the average income risk sharing coefficient was non-significant in the Sorensen et al (HIER LINK EINFÜGEN) paper 
in the original sample, we belive that the change may potentially come from data differences and the choice of base years.\\

Looking at table 5 we can see that the same results hold for the foreign assets over GDP regressions, giving us confidence to 
conclude that the negative relationship between Home Bias and Income and Consumption Risk sharing holds for the extended period 
of 30 years between 1987--2017.

% ============================================================
\subsection{Second Extention: Extended Country Sample 1987--2017}
% ============================================================
Further, we aim do analyse how the statistical relationship holds between 1987 and 2017 for a larger sample of 
OECD countries. With respect to the panel regressions, extending the sample period beyond 29 countries was not 
possible due to lacking data coverage either for GNI or for the Equity Home Bias measures. A starting point for
future research would defenitely be a larger extention for economies not included in our dataset, especially given the 
fact that the relative importance of non OECD economies has increased over the past two decades since the paper was published. 
The panel-regression results of the (slightly) extended sample between 1987 and 2017 confirm the same results as shown in the 
two regressions before. 
An decrease in the Equity Home Bias by 0.1 is associated with an increase of income risk sharing by roughly 3 percent and
an increase in consumption risk sharing by roughly 8 percent. 
again, the panel regressions for foreign Asset holdings relative to GDP confirm this estimation, with highly significant estimations 
for the asset holdings themselves, and non significant estimations for the time trend. Again, the average risk sharing coefficient for 
income is non significant in both specifications, while it is highly significant for consumption in both specifications.

% ============================================================
\section*{References}
% ============================================================

Add references manually here, or set up BibTeX/BibLaTeX later.

\end{document}